\textbf{Partie A}

\medskip
\begin{enumerate}
	\item Pour chaque caractère, il y a 64 possibilités, donc pour une séquences de 4 caractères, il y a $64^4$ possibilités, soit \np{16777216} possibilités. On pourra noter 
	
card$(\Omega)=\np{16777216}$.
	\item Si les caractères sont différents deux à deux, il s'agit alors d'un arrangement de 4 caractères parmi 64 : $\dfrac{64~!}{(64-4)~!}=64 \times 63 \times 62 \times 61=\np{15249024}$ possibilités.
	\item 
	\begin{enumerate}
		\item On reprend la question \textbf{1.} avec seulement 63 caractères. cela donne donc

$63^4=\np{15752961}$ possibilités.
		\item Soit $X$ l'évènement : \og la séquence ne comporte pas la lettre A \fg.

Son événement contraire est donc : $\overline{X}$ \og la séquence comporte au moins une fois la lettre A \fg{}.

		Ainsi card$(X) + \text{card}(\overline{X})= \text{card}(\Omega)$ donc card$(\overline{X}) = \text{card}(\Omega) - \text{card}(X) = 64^4 - 63^4 = \np{1024255}$.

Il y a donc \np{1024255} séquences contenant au moins une fois la lettre A.
		\item La lettre A peut se situer dans l'une des quatre positions dans le code, les trois autres lettres étant différentes. Il y aura donc $4\times 63^3 = \np{1000188}$ possibilités.
		\item Il y a $\dbinom{4}{2}$ façons de placer les deux lettres A dans la séquence. Les deux autres lettres ne sont pas des A. Il y aura donc $\dbinom{4}{2} \times 63^2 = \np{23814}$ possibilités.
	\end{enumerate}
\end{enumerate}

\medskip

\textbf{Partie B}

\medskip

\begin{enumerate}
	\item $X$ suit une loi binomiale de paramètres $n=250$ et $p=0,01$.
	\item $P(X=0)=\dbinom{250}{0}\times 0,01^0\times (1-0,01)^{250-0}=0,99^{250}\approx 0,081$.
	\item On calcule $P(X >16)= P(X\geqslant 17)$. Avec la calculatrice, $P(X\geqslant 17)\approx 1,04\times 10^{-9}$, ce qui est négligeable.
\end{enumerate}

\medskip

\textbf{Partie C}

\medskip

En utilisant les propriétés sur la linéarité de l'espérance et de la variance (les variables aléatoires $X_1,\: X_2,\: X_3$\: et $X_4$ étant indépendantes), on peut affirmer que :

$E(S)=E(X_1 + X_2 + X_3 + X_4)=E(X_1) + E(X_2) + E(X_3) + E(X_4)= 4\times E(X)=$

$ 4\times n \times p=4\times 250 \times 0,01 = 10$.

Il en est de même pour la variance :

$V(S)=V(X_1 + X_2 + X_3 + X_4)=V(X_1) + V(X_2) + V(X_3) + V(X_4)=4\times V(X)=4\times n \times p \times (1-p)=4\times 250 \times 0,01 \times 0,99=9,9$.
